\chapter{機率、條件機率與貝氏定理}
    在學會如何透過敘述統計和圖表來彙整資料之後，我們現在準備踏入推論統計的領域。推論統計的核心，在於從資料中汲取結論，同時以系統化的方式處理不確定性。因此，要進行推論統計，我們必須先熟悉量化不確定性的數學工具——機率。在本章中，我們將首先回顧機率與條件機率的定義，接著說明貝氏定理及其源自條件機率的實際應用。
    
    \begin{introduction}[第 \thechapter 章學習目標]{1}
        \item 區分機率的古典、頻率學派與主觀三種詮釋方式
        \item 理解事件機率的計算方式，以及如何定義事件的獨立和互斥
        \item 理解條件機率與貝氏定理，及其在醫學與流行病學中的應用
    \end{introduction}

\section{基礎機率}
\subsection{機率的定義}
    \textit{機率}(probability)又稱或然率或是概率，用來度量一個事件發生可能性，定義為介於 0 和 1 之間的實數。機率越接近 1，代表事件發生的可能性越高，而機率等於 1 代表事件必然發生；反之，機率越接近 0，代表事件發生的可能性越低，而機率等於 0 代表事件不可能發生。對於機率的解釋，一般可以分成兩大流派：\textit{客觀機率}(objective probability)和\textit{主觀機率}(subjective probability)。其中客觀機率又可以分成兩種定義方式：\textit{古典機率}(classical probability) 和\textit{頻率機率} (frequentist probability)。我們可以用「擲兩個骰子得到之總點數為 3 的機率」來說明這三種機率定義方式：
    \begin{itemize}\sloppy
        \item \textbf{古典機率}：古典機率是最原始的機率定義方法，此定義可以追溯到18世紀的數學家拉普拉斯(Pierre-Simon Laplace)。古典機率認為，首先要找到一群發生可能性均相同的簡單事件(simple events)，例如有 $n$ 個這樣的簡單事件。然後，我們可以將我們有興趣的事件分解為 $k$ 個簡單事件的組合。這樣一來，我們有興趣的事件的發生機率就是 $k/n$。例如，擲兩個骰子後，將第一個骰子的點數 a 和第二個骰子的點數 b 記為 (a,b)，就可以建構出 36 種可能性相同的簡單事件：(1,1), (1,2), (1,3), ... (1,6), (2,1), (2,2), ..., (6,6)。而我們有興趣的事件「得到總點數為 3」可以分解為兩種簡單事件 (1,2) 和 (2,1)，所以「擲兩個骰子得到之總點數為 3 的機率」等於 $2/36 = 1/18$。這種定義雖然簡單直觀，但有兩個主要缺點：首先，「可能性均相同的簡單事件」含有「可能性」在裡面，所以這種定義方法有自己指涉自己的套套邏輯之嫌；另外，符合定義的簡單事件不見得存在，例如在上述的例子中，如果骰子是灌鉛的骰子，那麼就找不到可能性相同的簡單事件了。
        \item \textbf{頻率機率}：頻率機率認為，要探討事件的機率，首先必須要有可重複的隨機試驗，例如我們這裡的隨機試驗就是丟兩顆骰子並記錄點數。該隨機試驗的所有可能結果所組成的集合叫做\textit{樣本空間}(sample space)，通常記為$\Omega$，因此例子中的樣本空間就是 36 種可能點數組合的集合，也就是 $\Omega = \{(1,1), (1,2), (1,3), ... (1,6), (2,1), (2,2), ..., (6,6)\}$。注意到這裡我們並未要求樣本空間中所有結果的可能性相等，因此和古典機率不同。\textit{事件}(event)就是樣本空間的子集合，例子中「得到總點數為 3」的事件即對應到 $A = \{(1,2),(2,1)\}$ 這個子集合。重複該試驗多次以後，事件發生的次數 $k$ 和試驗次數 $n$ 的比值雖可能不固定，但是隨著試驗次數增多終將會趨近一個數值，而該數值就是事件發生的機率。頻率機率的定義解決了古典機率的兩大缺點，但其奠基於「可重複隨機試驗」的事實也引起批評：如果我們關心的事件不可能重複試驗，例如「明天台北發生地震」、「某國總統被成功彈劾」或「某病患將在 24 小時內離開加護病房」，那麼頻率機率的定義就顯得不合適。
        \item \textbf{主觀機率}：主觀機率認為機率是個人信念的一種度量，因此「擲兩個骰子得到之總點數為 3 的機率」是我們基於已知的事實（例如骰子的製作過程、投擲的隨機性）評估出擲出骰子後總點數為 3 的可能性。和頻率機率的差異在於，頻率機率認為存在可重複的隨機試驗，該試驗的條件可能未知，但已固定。主觀機率則認為針對試驗的方方面面的不確定性和主觀認知，都應該納入機率考量。因此不同人針對同一事件的主觀機率，也會隨著每個人的信念和經驗而不同。
    \end{itemize}
    在後續的課程中，大部分時候我們選擇的機率定義為頻率機率，但必要的時候，我們也會改用主觀機率以方便實務上的解釋。讀者閱讀時也不妨想想各處提到的機率是用哪種方式定義比較合適。
\subsection{機率的基本性質}
    在頻率機率的定義下，假設樣本空間是 $\Omega$，而 $A$, $B$ 代表任意的事件，我們定義如下：
    \begin{itemize}
        \item $\PP(A)$ 代表 $A$ 事件發生的機率。
        \item $\bar{A} = \Omega \backslash A$ 為 $A$ 的\textit{補事件} (complementary event)，也就是 「$A$沒有發生」的事件。
        \item $A \cap B$ 為 $A$ 和 $B$ 兩個事件的交集，對應到「$A$與$B$均發生」的事件。
        \item $A \cup B$ 為 $A$ 和 $B$ 兩個事件的聯集，對應到「$A$或$B$發生」的事件。
        \item 當 $A \cap B = \emptyset$，即兩個事件的交集為空集合時，我們稱 $A$ 和 $B$ 兩個事件\textit{互斥} (mutually exclusive)。此時這兩個事件不可能同時發生，因為他們代表的結果完全沒有交集。
    \end{itemize}
    舉例來說，如果 $\Omega = \{1,2,3,4,5,6\}$ 是擲一顆骰子的點數樣本空間，$A = \{2,4,6\}$ 是擲到偶數的事件，$B = \{1,2,3\}$ 是擲到小點的事件，$C = \{5\}$ 是擲到 5 點的事件。那麼 $A$ 的補事件是 $\bar{A} = \{1,3,5\}$，也就是擲到奇數的事件。$A \cap B = \{2\}$ 是擲到偶數小點的事件。$A \cup B = \{1,2,3,4,6\}$ 是擲到偶數或小點的事件。而因為 $A \cap C = \emptyset$，所以 $A$ 和 $C$ 是互斥事件。

    根據上述的定義，下面幾個機率性質必須要成立：
    
    \begin{itemize}
        \item $\PP(\Omega) = 1$：樣本空間本身就是樣本空間的子集合，所以樣本空間也是一個事件（也就是觀察到所有可能結果的組合）。這個性質說明\textbf{所有可能結果的總機率應為 1} 。
        \item $0 \le \PP(A) \le 1$：這個性質說明\textbf{機率必然是0到1之間的實數}。
        \item $A \cap B = \phi \Rightarrow \PP(A \cup B) = \PP(A) + \PP(B)$：這個性質說明\textbf{兩個事件若互斥，則至少一個事件發生的機率為兩個事件的機率相加}。
    \end{itemize}

    根據這些性質，因為 $A \cap \bar{A} = \phi$，所以 $\PP(A) + \PP(\bar{A}) = \PP(A \cup \bar{A}) = \PP(\Omega) = 1$，也就是不發生 $A$ 事件的機率和發生 $A$ 事件的機率總和為 1，非常符合直覺。另外，根據簡單的集合性質我們也能證明，在 $A$ 與 $B$ 不互斥的情況下的\textit{加法原則} (addition rule)：
    
    \[\PP(A \cup B) = \PP(A) + \PP(B) - \PP(A \cap B)\]


\section{條件機率與獨立事件}

    接下來我們要說明一個比較複雜的概念：\textit{條件機率} (conditional probability)。條件機率評估的是「給定某事件發生的情況下，另外一個事件發生的機率」，例如「在 COVID-19 快篩呈陽性的情況下，真的感染 COVID-19 的機率」。以數學符號來表示的話，給定 $B$ 事件下 $A$ 事件發生的條件機率記為 $\PP(A|B)$，其中「$A|B$」念為「A 給定 B（A given B）」。條件機率的計算方法十分直覺：假設我們重複做實驗做了 $n$ 次，其中大約會有 $n\PP(B)$ 次發生 $B$ 事件，而約有 $n\PP(A \cap B)$ 次除了發生 $B$ 事件、還發生了 $A$ 事件。因此，給定 $B$ 事件下 $A$ 事件發生的機率就可以寫為
    
    \[\PP(A|B) = \frac{n\PP(A \cap B)}{n\PP(B)} = \frac{\PP(A \cap B)}{\PP(B)}\]

    \noindent 舉例來說，如果我們想知道擲兩枚公正骰子，給定點數和為 4 的情況下，兩個骰子的點數不同的機率為多少，也就是$\PP(\text{兩個骰子點數不同}|\text{點數和為4})$。簡單計算可以得知，點數和為 4 的機率 $\PP(\text{點數和為4}) = 3/36$。點數和為 4 且兩個骰子點數不同的機率為 $\PP(\text{(兩個骰子點數不同)}\cap\text{(點數和為4)}) = 2/36$。所以
    
    \[\PP(\text{兩個骰子點數不同}|\text{點數和為4}) = \frac{\PP(\text{(兩個骰子點數不同)}\cap\text{(點數和為4)})}{\PP(\text{點數和為4})} = \frac{2}{3}\]

    \bigskip

    \begin{custom}{動動腦}
        按照條件機率的定義，如果 $\PP(B) \ne 0$，$\PP(A|B)$ 是否可能大於 1？為什麼？
    \end{custom}

    \answerbox{120pt}

    \bigskip

    如果一個事件發生與否，並不會影響到另一個事件發生的機率，那麼這兩個事件應該互不影響。從這個角度來看，兩個事件 $A$ 和 $B$ 互不影響應該意謂著 $\PP(A|B) = \PP(A)$ 而且 $\PP(B|A) = \PP(B)$。根據條件機率的定義，這兩個式子都導向同一個結果：

    \[\PP(A \cap B) = \PP(A) \PP(B)\]

    \noindent 如果兩個事件 $A$ 和 $B$ 滿足上述關係，他們就被稱為\textit{獨立} (independent)事件。反之，則被稱為\textit{不獨立}或\textit{相依} (dependent)事件。在前面的例子中，兩個骰子點數不同的機率為 $30/36$，點數和為 4 的機率為 $3/36$，但骰子點數不同且點數和為 4 的機率為 $2/36$。顯然 $(30/36)(3/36) \ne 2/36$，所以「兩個骰子點數不同」和「點數和為 4」這兩個事件不獨立。
    
    \bigskip

    \begin{custom}{動動腦}
        互斥事件是否有可能是獨立事件？如果有可能，是在什麼情況下才會成立？
    \end{custom}

    \answerbox{120pt}

    \bigskip

    條件機率在流行病學的應用十分重要，尤其是篩檢診斷結果的判讀。臨床上有許多疾病的黃金診斷標準十分耗費資源，因此就公共衛生角度而言，較佳的策略是先進行一個準確度較低的篩檢，待篩檢陽性後再進行較昂貴耗時的診斷檢查。例如 HIV 的黃金診斷標準需經過抗體免疫層析或分子生物學核酸檢測，但有疑慮的感染者則建議先採指尖血或唾液做初步篩檢，呈現陽性後再進行確認檢驗。篩檢結果的結果判讀中，有下列的重要參數需要評估（以表\ref{tab:screening_high_risk}為例）：
    \begin{itemize}\sloppy
        \item \textit{盛行率} (Prevalence)：$\PP(\text{有疾病})$，即總體受檢者有疾病的比例。在下表中為 $7000/10000 = 0.7$。
        \item \textit{敏感度} (Sensitivity)：$\PP(\text{篩檢陽性}|\text{有疾病})$，即有疾病的人篩檢呈陽性的機率。在下表中為 $5600/7000 = 0.8$。
        \item \textit{特異度} (Specificity)：$\PP(\text{篩檢陰性}|\text{無疾病})$，即無疾病的人篩檢呈陰性的機率。在下表中為 $2700/3000 = 0.9$。
        \item \textit{陽性預測值} (Positive predictive value, PPV)：$\PP(\text{有疾病}|\text{篩檢陽性})$，即篩檢陽性的人有疾病的機率。在下表中為 $5600/5900 \approx 0.949$。
        \item \textit{陰性預測值} (Negative predictive value, NPV)：$\PP(\text{無疾病}|\text{篩檢陰性})$，即篩檢陰性的人無疾病的機率。在下表中為 $2700/4100 \approx 0.659$。
    \end{itemize}

    \begin{table}[htbp]
        \begin{center}
            \begin{tabular}{c|cc|c}
                \toprule
                 & \textbf{有疾病} & \textbf{無疾病} & 總和\\
                \hline
                \textbf{篩檢陽性} & 5600 & 300 & 5900\\
                \textbf{篩檢陰性} & 1400 & 2700 & 4100\\
                \hline
                總和 & 7000 & 3000 & 10000\\
                \bottomrule
            \end{tabular}
            \caption{高風險族群篩檢結果\label{tab:screening_high_risk}}
        \end{center}
    \end{table}

    從這裡可以觀察到，$\PP(A|B)$ 和 $\PP(B|A)$ 是兩個完全不同的機率。此處篩檢陽性時患病的機率很高，到達約 0.95，但即便篩檢工具的特異度已經高達 0.9，但是在盛行率較高的高風險族群中，篩檢陰性者的沒病機率仍僅有約 0.66。反之，如果我們使用同樣敏感度和特異度篩檢工具，但用在風險較低的族群中，結果可能會不一樣。下面我們實作看看：

    \bigskip
    
    \begin{custom}{動手作}
        假設我們同樣有 10000 人的族群，但盛行率降為 0.3，且使用表 \ref{tab:screening_high_risk} 的篩檢工具進行篩檢（敏感度與特異度不變），請填入下面表 \ref{tab:screening_low_risk} 各空格並計算 PPV 與 NPV。
    \end{custom}

    \begin{table}[htbp]
        \begin{center}
            \begin{tabular}{c|cc|c}
                \toprule
                 & \textbf{有疾病} & \textbf{無疾病} & 總和\\
                \hline
                \textbf{篩檢陽性} & \phantom{2400} & \phantom{700} & \phantom{3100}\\
                \textbf{篩檢陰性} & \phantom{600} & \phantom{6300} & \phantom{6900}\\
                \hline
                總和 & \phantom{3000} & \phantom{7000} & 10000\\
                \bottomrule
            \end{tabular}
            \caption{低風險族群篩檢結果\label{tab:screening_low_risk}}
        \end{center}
    \end{table}

    可以看到，盛行率下降的情況下，陰性預測值陡然攀升至 0.913，而陽性預測值則變得不如預期。因此，判讀篩檢結果時，除了篩檢工具的敏感度和特異度外，還要將盛行率納入考量。舉例來說，如果在疫情和緩而且完全無症狀的情況下進行 COVID 快篩，此時盛行率相對很低，因此即使篩檢陽性，實際上感染 COVID 的機會也很低；反之，如果疫情正在快速蔓延且症狀明確，則盛行率相對較高，當篩檢陽性時就很可能已經感染。

    \bigskip

    \begin{custom}{動動腦}
        雖然理想狀態下，我們希望篩檢工具同時具有高敏感度及高特異度，但實際上在篩檢工具的設計上，這兩者通常是個權衡關係。如果一個疾病的黃金標準檢驗很昂貴，但有病時並不需要太急著治療，那麼我們設計篩檢的目的應該要儘量確定陰性者沒病（提高 NPV）還是確定陽性者有病（提高 PPV），此時應該選擇提高篩檢工具的敏感度還是特異度？反之，如果一個疾病的黃金標準檢驗相對便宜，但有病時必須要馬上給予治療，此時的決策應為何？
    \end{custom}

    \answerbox{240pt}

\section{貝氏定理}
    在前面的例子中，我們為了算出盛行率降低後陽性預測值和陰性預測值的變化，藉由假設總族群有 10000 人，算出表格中的所有人數後，再依據定義算出陽性和陰性預測值。如果我們把這些數字都符號化，將總人數記作 $n$，將有疾病的事件寫作 $D$，篩檢陽性的事件寫作 $T$，那麼表格就會變成如表 \ref{tab:Bayes}：

    \begin{table}[htbp]
        \begin{center}
            \begin{tabular}{c|cc|c}
                \toprule
                 & $D$ & $\bar{D}$ & 總和\\
                \hline
                $T$ & $n\PP(D)\PP(T|D)$ & $n\PP(\bar{D})\PP(T|\bar{D})$ & $n[\PP(D)\PP(T|D) + \PP(\bar{D})\PP(T|\bar{D})]$\\
                $\bar{T}$ & $n\PP(D)\PP(\bar{T}|D)$  & $n\PP(\bar{D})\PP(\bar{T}|\bar{D})$ & $n[\PP(D)\PP(\bar{T}|D) + \PP(\bar{D})\PP(\bar{T}|\bar{D})]$\\
                \hline
                總和 & $n\PP(D)$ & $n\PP(\bar{D})$ & $n$\\
                \bottomrule
            \end{tabular}
            \caption{貝氏定理的說明\label{tab:Bayes}}
        \end{center}
    \end{table}
    
    由表中的數據試著計算陽性預測值，就可以得到著名的\textit{貝氏定理}(Bayes' Theorem)：
    
    \[\PP(D|T) = \frac{\PP(T|D)\PP(D)}{\PP(T|D)\PP(D) + \PP(T|\bar{D})\PP(\bar{D})}\]
    
    \noindent 貝氏定理最早是由英國的牧師兼數學家貝葉斯（Thomas Bayes, 1701-1761）提出，並由他的友人替他整理發表。在該定理中，我們可以把分子的 $\PP(D)$ 移出來，把定理寫成
    
    \[\PP(D|T) = \Bigg(\frac{\PP(T|D)}{\PP(T|D)\PP(D) + \PP(T|\bar{D})\PP(\bar{D})}\Bigg) \PP(D)\]
    
    \noindent 在我們做篩檢前，對於族群的每個人的患病機率只能用總體的盛行率 $\PP(D)$ 估計，因此 $\PP(D)$ 也被稱為患病的\textit{先驗機率} (prior probability)。做篩檢後得知病患篩檢陽性後，我們對該病患的患病機率估計變成了 $\PP(D|T)$，因此 $\PP(D|T)$ 也被稱為得知篩檢陽性後患病的\textit{後驗機率} (posterior probability)。而中間大括弧的部分就是將先驗機率\textit{更新}(update)為後驗機率的因子，也可看作篩檢陽性帶給我們的資訊。例如在表\ref{tab:screening_high_risk}中，各病患的患病先驗機率是 0.7，一旦篩檢呈現陽性，患病的後驗機率就被更新到 0.949。綜上所述，我們可以寫作如下：
    
    \[\text{後驗患病機率}=\text{篩檢陽性的更新因子}\times\text{先驗患病機率}\]
    
    \noindent 不過注意到，這裡「篩檢陽性的更新因子」和盛行率 $\PP(D)$ 有關，因此同一個篩檢工具在不同族群的更新因子會有所不同。換言之，這個更新因子無法代表篩檢工具的核心特性。為了解決這個問題，我們寫出貝氏定理針對未患病機率的公式（把 $D$ 換成 $\bar{D}$），可得到：
    
    \[\PP(\bar{D}|T) = \Bigg(\frac{\PP(T|\bar{D})}{\PP(T|\bar{D})\PP(\bar{D}) + \PP(T|D)\PP(D)}\Bigg) \PP(\bar{D})\]
    
    \noindent 將這個式子與前面的式子相除後，即得：
    
    \[\frac{\PP(D|T)}{\PP(\bar{D}|T)} = \frac{\PP(T|D)}{\PP(T|\bar{D})} \times \frac{\PP(D)}{\PP(\bar{D})}\]
    
    \noindent 等式右邊第二項 $\frac{\PP(D)}{\PP(\bar{D})}$ 是先驗的患病機率除以未患病機率。一般我們把「發生機率除以未發生機率」的形式稱做\textit{勝算}(odds)，例如一枚公正銅板出現正面的勝算為 1，一顆公正骰子擲出一點的勝算為 1/5。因此，$\frac{\PP(D)}{\PP(\bar{D})}$ 是先驗的患病勝算。同理，等號左邊的 $\frac{\PP(D|T)}{\PP(\bar{D}|T)}$ 則是後驗的患病勝算。中間的 $\frac{\PP(T|D)}{\PP(T|\bar{D})}$中，分子是「患病者篩檢陽性的機率（敏感度）」，分母則是「未患病者篩檢陽性的機率（一減特異度）」，所以整個分數代表「患病者相對於未患病者的篩檢陽性機率比值」。這個分數不再與盛行率相關，而僅和篩檢工具的敏感度和特異度有關，因此可視為篩檢工具呈現陽性時特有的診斷力特徵，被稱為\textit{陽性概似比} (positive likelihood ratio, LR+)。綜上所述，在篩檢陽性的時候我們有：
    
    \[\text{後驗患病勝算}=\text{陽性概似比(LR+)}\times\text{先驗患病勝算}\]
    
    \noindent 同理，在篩檢陰性的時候我們有
    
    \[\frac{\PP(D|\bar{T})}{\PP(\bar{D}|\bar{T})} = \frac{\PP(\bar{T}|D)}{\PP(\bar{T}|\bar{D})} \times \frac{\PP(D)}{\PP(\bar{D})}\]
    
    \noindent 此處我們稱$\frac{\PP(\bar{T}|D)}{\PP(\bar{T}|\bar{D})}$，也就是「一減敏感度」除以「特異度」為陰性概似比(negative likelihood ratio, LR-)，因此在篩檢陰性的時候我們可寫作
    
    \[\text{後驗患病勝算}=\text{陰性概似比(LR-)}\times\text{先驗患病勝算}\]
    
    \noindent 由此可見，只要知道篩檢工具的陽性及陰性概似比，就可以根據病患的先驗患病勝算及檢驗結果推斷出其後驗患病勝算。因此實證醫學中，這兩個數值是評估篩檢或診斷工具有效度的重要指標。有關貝氏定理的更多延伸，請參考以下\href{https://www.youtube.com/watch?v=R13BD8qKeTg}{影片}。

    \bigskip

    \begin{custom}{動動腦}
        陽性概似比和陰性概似比的可能取值範圍為何？針對一個篩檢工具，我們會希望它的陽性概似比和陰性概似比越大越好、還是越小越好？如果有一個篩檢工具的陽性概似比為 1、陰性概似比為 0.005，我們應該要如何使用這個篩檢工具的檢驗結果？
    \end{custom}

    \answerbox{180pt}

    \begin{custom}{動動腦}
        Let's make a deal 是美國自 1960 年代開播的長壽遊戲節目。節目的第一代主持人是 Monty Hall，至今已經播出超過六千集。其中一個遊戲環節中，節目組會在左、中、右三扇門後面隨機地放兩隻山羊和一輛跑車。參賽者需要選定一扇門，以獲得門後的獎品。（假定參賽者想要跑車而不是山羊！）參賽者選定門後，Monty 會在剩下的兩扇門中開啟一扇，顯示出後面的山羊，並問參賽者是不是要換門。你做為參賽者選了左邊的門，而後 Monty 開了中間的門顯示出一隻山羊：
        
        (A) [經典 Monty] 如果 Monty 總是會開門顯示山羊，而且剩下的兩扇門都是山羊的話就隨便開一扇，那麼你應不應該換門？
        
        (B) [左撇子 Monty] 如果 Monty 總是會開門顯示山羊，而且剩下的兩扇門都是山羊的話就開靠左的那扇，那麼你應不應該換門？

        (C) [無知 Monty] 如果 Monty 根本不知道門後是什麼，只總是隨便亂開一扇門（只是剛好節目中都出現山羊），那麼你應不應該換門？
        
        (D) [騙子 Monty] 如果 Monty Hall 在參賽者選對了都給換門（門後兩頭羊時隨機開門），選錯就不給機會，那麼你應不應該換門？
                
        (E) [隨機 Monty] 如果 Monty Hall 在參賽者選對時有 $p_1$ 的機率給換門（而且隨機開門），選錯時有 $p_2$ 的機率給換門，你應該如何根據 $p_1$ 和 $p_2$ 選擇是否換門？

        (F) 承上題，如果 Monty 要讓節目送出車子的機率最小化，應該如何設定 $p_1$ 和 $p_2$？

        (G) 承上題，如果節目組希望 Monty 能在最小化送車機率的同時，盡量多玩換門遊戲以增加節目可看性，那麼 Monty 應該如何設定 $p_1$ 和 $p_2$？

    \end{custom}

    \answerbox{320pt}
    
    \bigskip

    \begin{takehome}
        \item 機率的定義主要有三種方式：古典機率、頻率機率與主觀機率。
        \item 聯集事件的加法原則：$\PP(A \cup B) = \PP(A) + \PP(B) - \PP(A \cap B)$
        \item 「兩個事件 $A$ 與 $B$ 互相獨立」等價於 $\PP(A \cap B) = \PP(A)\PP(B)$
        \item 條件機率是在另一事件已發生的條件下，某事件發生的機率，其定義為 $\PP(A|B) = \PP(A \cap B) / \PP(B)$
        \item 貝氏定理說明先驗機率（信念）可藉由觀測到的資訊更新為後驗機率（信念）。
        \item 評估診斷或篩檢工具表現與結果的重要參數：敏感度、特異度、陽性概似比、陰性概似比、陽性預測值、陰性預測值。
    \end{takehome}

    \begin{docexam}{（104-1醫學(一)）}
        若已知國內男性抽菸盛行率為 $20\%$，隨機抽三位男性，三位都抽菸的機率為何？
    \end{docexam}

    \begin{docexam}{（102-1醫學(一)）}
        某學生想為他這學期體育課的選課做出決定。假設一學期可選兩種運動課程，但選課人數一經額滿就無法選上。若此生選上游泳課的機率為 0.6，選上韻律課的機率為 0.5，同時選上游泳課和韻律課的機率為 0.3。此生選上韻律課或游泳課或兩種課程的機率為多少？
    \end{docexam}

    \begin{docexam}{（109-2醫學(二)）}
        假設有一個疾病為顯性遺傳病，得病者父母親中，其中有一位有病，另一位則無。這種疾病每位子女遺傳的機率 $1/2$，一個有 2 個小孩的家庭，兩個小孩皆得病的機率為何？
    \end{docexam}